\Chapter{1}

\Verse{1}
{
    \zA{――此开卷第一回也。 作者自云：曾历过一番梦幻之后，故将真事隐去，而借 通灵说此《石头记》一书也，故曰“甄士隐”云云。 但书中所记何事何人？ 自己又
        云：“今风尘碌碌，一事无成，忽念及当日所有之女子：一一细考较去，觉其行止 见识皆出我之上。 我堂堂须眉诚不若彼裙钗，我实愧则有馀，悔又无益，大无可如
        何之日也。 当此日，欲将已往所赖天恩祖德，锦衣纨？}
}
{
    \eA{This is the opening section; this the first chapter. Subsequent to the
        visions of a dream which he had, on some previous occasion, experienced, the
        writer personally relates, he designedly concealed the true circumstances,
        and borrowed the attributes of perception and spirituality to relate this
        story of the Record of the Stone.}
}
{
}

\Verse{2}
{
    \zA{之时，饫甘餍肥之日，背父 兄教育之恩，负师友规训之德，以致今日一技无成、半生潦倒之罪，编述一集，以 告天下；
        知我之负罪固多，然闺阁中历历有人，万不可因我之不肖，自护己短，一 并使其泯灭也。 所以蓬牖茅椽，绳床瓦灶，并不足妨我襟怀； 况那晨风夕月，阶柳
        庭花，更觉得润人笔墨。 我虽不学无文，又何妨用假语村言敷演出来？ 亦可使闺阁 昭传。 复可破一时之闷，醒同人之目，不亦宜乎？” 故曰“贾雨村”云云。}
}
{
    \eA{With this purpose, he made use of such designations as Chen Shih-yin (truth
        under the garb of fiction) and the like. What are, however, the events
        recorded in this work? Who are the dramatis personae?}
}
{
}

\Verse{3}
{
    \zA{更于篇 中间用“梦”“幻”等字，却是此书本旨，兼寓提醒阅者之意。 看官你道此书从何而起？ 说来虽近荒唐，细玩颇有趣味。 却说那女娲氏炼石补
        天之时，于大荒山无稽崖炼成高十二丈、见方二十四丈大的顽石三万六千五百零一 块。 那娲皇只用了三万六千五百块，单单剩下一块未用，弃在青埂峰下。 谁知此石
        自经锻炼之后，灵性已通，自去自来，可大可小。 因见众石俱得补天，独自己无才 不得入选，遂自怨自愧，日夜悲哀。}
}
{
    \eA{Wearied with the drudgery experienced of late in the world, the author
        speaking for himself, goes on to explain, with the lack of success which
        attended every single concern, I suddenly bethought myself of the womankind
        of past ages. Passing one by one under a minute scrutiny, I felt that in
        action and in lore, one and all were far above me;}
}
{
}

\Verse{4}
{
    \zA{一日正当嗟悼之际，俄见一僧一道远远而来， 生得骨格不凡，丰神迥异，来到这青埂峰下，席地坐谈。 见着这块鲜莹明洁的石头， 且又缩成扇坠一般，甚属可爱。
        那僧托于掌上，笑道：“形体倒也是个灵物了，只 是没有实在的好处。 须得再镌上几个字，使人人见了便知你是件奇物，然后携你到
        那昌明隆盛之邦、诗礼簪缨之族、花柳繁华地、温柔富贵乡那里去走一遭。” 石头 听了大喜，因问：“不知可镌何字？ 携到何方？ 望乞明示。”}
}
{
    \eA{that in spite of the majesty of my manliness, I could not, in point of fact,
        compare with these characters of the gentle sex. And my shame forsooth then
        knew no bounds; while regret, on the other hand, was of no avail, as there
        was not even a remote possibility of a day of remedy.}
}
{
}

\Verse{5}
{
    \zA{那僧笑道：“你且莫问， 日后自然明白。” 说毕，便袖了，同那道人飘然而去，竟不知投向何方。
        又不知过了几世几劫，因有个空空道人访道求仙，从这大荒山无稽崖青埂峰下 经过。 忽见一块大石，上面字迹分明，编述历历。 空空道人乃从头一看，原来是无
        才补天、幻形入世，被那茫茫大士、渺渺真人携入红尘、引登彼岸的一块顽石； 上 面叙着堕落之乡、投胎之处，以及家庭琐事、闺阁闲情、诗词谜语，倒还全备。 只
        是朝代年纪，失落无考。}
}
{
    \eA{On this very day it was that I became desirous to compile, in a connected
        form, for publication throughout the world, with a view to (universal)
        information, how that I bear inexorable and manifold retribution;}
}
{
}

\Verse{6}
{
    \zA{后面又有一偈云： 无才可去补苍天，枉入红尘若许年。 此系身前身后事，倩谁记去作奇传？
        空空道人看了一回，晓得这石头有些来历，遂向石头说道：“石兄，你这一段故事， 据你自己说来，有些趣味，故镌写在此，意欲闻世传奇。 据我看来：第一件，无朝
        代年纪可考； 第二件，并无大贤大忠、理朝廷、治风俗的善政，其中只不过几个异 样女子，或情或痴，或小才微善。 我纵然抄去，也算不得一种奇书。”}
}
{
    \eA{inasmuch as what time, by the sustenance of the benevolence of Heaven, and
        the virtue of my ancestors, my apparel was rich and fine, and as what days
        my fare was savory and sumptuous, I disregarded the bounty of education and
        nurture of father and mother, and paid no heed to the virtue of precept and
        injunction of teachers and friends, with the result that I incurred the
        punishment, of failure recently in the least trifle, and the reckless waste
        of half my lifetime.}
}
{
}

\Verse{7}
{
    \zA{石头果然答 道：“我师何必太痴!我想历来野史的朝代，无非假借汉、唐的名色； 莫如我这石头 所记不借此套，只按自己的事体情理，反倒新鲜别致。
        况且那野史中，或讪谤君相， 或贬人妻女，奸淫凶恶，不可胜数； 更有一种风月笔墨，其淫秽污臭最易坏人子弟。
        至于才子佳人等书，则又开口‘文君’，满篇‘子建’，千部一腔，千人一面，且终 不能不涉淫滥。 在作者不过要写出自己的两首情诗艳赋来，故假捏出男女二人名姓；}
}
{
    \eA{There have been meanwhile, generation after generation, those in the inner
        chambers, the whole mass of whom could not, on any account, be, through my
        influence, allowed to fall into extinction, in order that I, unfilial as I
        have been, may have the means to screen my own shortcomings.}
}
{
}

\Verse{8}
{
    \zA{又必旁添一小人拨乱其间，如戏中的小丑一般。 更可厌者，‘之乎者也’，非理即文， 大不近情，自相矛盾。
        竟不如我这半世亲见亲闻的几个女子，虽不敢说强似前代书 中所有之人，但观其事迹原委，亦可消愁破闷； 至于几首歪诗，也可以喷饭供酒。
        其间离合悲欢，兴衰际遇，俱是按迹循踪，不敢稍加穿凿，至失其真。 只愿世人当 那醉馀睡醒之时，或避事消愁之际，把此一玩，不但是洗旧翻新，却也省了些寿命
        筋力，不更去谋虚逐妄了。 我师意为如何？”}
}
{
    \eA{Hence it is that the thatched shed, with bamboo mat windows, the bed of tow
        and the stove of brick, which are at present my share, are not sufficient to
        deter me from carrying out the fixed purpose of my mind. And could I,
        furthermore, confront the morning breeze, the evening moon, the willows by
        the steps and the flowers in the courtyard, methinks these would moisten to
        a greater degree my mortal pen with ink;}
}
{
}

\Verse{9}
{
    \zA{空空道人听如此说，思忖半晌，将这《石头记》再检阅一遍。 因见上面大旨不 过谈情，亦只是实录其事，绝无伤时诲淫之病，方从头至尾抄写回来，闻世传奇。
        从此空空道人因空见色，由色生情，传情入色，自色悟空，遂改名情僧，改《石头 记》为《情僧录》。 东鲁孔梅溪题曰《风月宝鉴》。 后因曹雪芹于悼红轩中，披阅十
        载，增删五次，纂成目录，分出章回，又题曰《金陵十二钗》，并题一绝。 即此便 是《石头记》的缘起。 诗云： 满纸荒唐言，一把辛酸泪。}
}
{
    \eA{but though I lack culture and erudition, what harm is there, however, in
        employing fiction and unrecondite language to give utterance to the merits
        of these characters? And were I also able to induce the inmates of the inner
        chamber to understand and diffuse them, could I besides break the weariness
        of even so much as a single moment, or could I open the eyes of my
        contemporaries, will it not forsooth prove a boon?}
}
{
}

\Verse{10}
{
    \zA{都云作者痴，谁解其中味! 《石头记》缘起既明，正不知那石头上面记着何人何事？ 看官请听。 按那石上
        书云：当日地陷东南，这东南有个姑苏城，城中阊门，最是红尘中一二等富贵风流 之地。 这阊门外有个十里街，街内有个仁清巷，巷内有个古庙，因地方狭窄，人皆
        呼作“葫芦庙”。 庙旁住着一家乡宦，姓甄名费字士隐，嫡妻封氏，性情贤淑，深 明礼义。 家中虽不甚富贵，然本地也推他为望族了。}
}
{
    \eA{This consideration has led to the usage of such names as Chia Yü-ts'un and
        other similar appellations. More than any in these pages have been employed
        such words as dreams and visions; but these dreams constitute the main
        argument of this work, and combine, furthermore, the design of giving a word
        of warning to my readers. Reader, can you suggest whence the story begins?}
}
{
}

\Verse{11}
{
    \zA{因这甄士隐禀性恬淡，不以功 名为念，每日只以观花种竹、酌酒吟诗为乐，倒是神仙一流人物。 只是一件不足： 年过半百，膝下无儿，只有一女乳名英莲，年方三岁。
        一日炎夏永昼，士隐于书房闲坐，手倦抛书，伏几盹睡，不觉朦胧中走至一处， 不辨是何地方。 忽见那厢来了一僧一道，且行且谈。 只听道人问道：“你携了此物，
        意欲何往？” 那僧笑道：“你放心，如今现有一段风流公案正该了结，这一干风流 冤家尚未投胎入世。}
}
{
    \eA{The narration may border on the limits of incoherency and triviality, but it
        possesses considerable zest. But to begin. The Empress Nü Wo, (the goddess
        of works,) in fashioning blocks of stones, for the repair of the heavens,
        prepared, at the Ta Huang Hills and Wu Ch'i cave, 36,501 blocks of rough
        stone, each twelve chang in height, and twenty-four chang square. Of these
        stones, the Empress Wo only used 36,500;}
}
{
}

\Verse{12}
{
    \zA{趁此机会，就将此物夹带于中，使他去经历经历。” 那道人道： “原来近日风流冤家又将造劫历世，但不知起于何处，落于何方？” 那僧道：“此 事说来好笑。
        只因当年这个石头，娲皇未用，自己却也落得逍遥自在，各处去游玩。 一日来到警幻仙子处，那仙子知他有些来历，因留他在赤霞宫中，名他为赤霞宫神 瑛侍者。
        他却常在西方灵河岸上行走，看见那灵河岸上三生石畔有棵绛珠仙草，十 分娇娜可爱，遂日以甘露灌溉，这绛珠草始得久延岁月。}
}
{
    \eA{so that one single block remained over and above, without being turned to
        any account. This was cast down the Ch'ing Keng peak. This stone, strange to
        say, after having undergone a process of refinement, attained a nature of
        efficiency, and could, by its innate powers, set itself into motion and was
        able to expand and to contract.}
}
{
}

\Verse{13}
{
    \zA{后来既受天地精华，复得 甘露滋养，遂脱了草木之胎，幻化人形，仅仅修成女体，终日游于离恨天外，饥餐 秘情果，渴饮灌愁水。
        只因尚未酬报灌溉之德，故甚至五内郁结着一段缠绵不尽之 意。 常说：‘自己受了他雨露之惠，我并无此水可还。 他若下世为人，我也同去走
        一遭，但把我一生所有的眼泪还他，也还得过了。’ 因此一事，就勾出多少风流冤 家都要下凡，造历幻缘，那绛珠仙草也在其中。}
}
{
    \eA{When it became aware that the whole number of blocks had been made use of to
        repair the heavens, that it alone had been destitute of the necessary
        properties and had been unfit to attain selection, it forthwith felt within
        itself vexation and shame, and day and night, it gave way to anguish and
        sorrow.}
}
{
}

\Verse{14}
{
    \zA{今日这石正该下世，我来特地将他 仍带到警幻仙子案前，给他挂了号，同这些情鬼下凡，一了此案。” 那道人道：“果 是好笑，从来不闻有‘还泪’之说。
        趁此你我何不也下世度脱几个，岂不是一场功 德？” 那僧道：“正合吾意。 你且同我到警幻仙子宫中将这蠢物交割清楚，待这一 干风流孽鬼下世，你我再去。
        如今有一半落尘，然犹未全集。” 道人道：“既如此， 便随你去来。” 却说甄士隐俱听得明白，遂不禁上前施礼，笑问道：“二位仙师请了。”}
}
{
    \eA{One day, while it lamented its lot, it suddenly caught sight, at a great
        distance, of a Buddhist bonze and of a Taoist priest coming towards that
        direction. Their appearance was uncommon, their easy manner remarkable. When
        they drew near this Ch'ing Keng peak, they sat on the ground to rest, and
        began to converse.}
}
{
}

\Verse{15}
{
    \zA{那僧道 也忙答礼相问。 士隐因说道：“适闻仙师所谈因果，实人世罕闻者，但弟子愚拙， 不能洞悉明白。
        若蒙大开痴顽，备细一闻，弟子洗耳谛听，稍能警省，亦可免沉沦 之苦了。” 二仙笑道：“此乃玄机，不可预泄。 到那时只不要忘了我二人，便可跳出 火坑矣。”
        士隐听了，不便再问，因笑道：“玄机固不可泄露，但适云‘蠢物’，不 知为何，或可得见否？” 那僧说：“若问此物，倒有一面之缘。” 说着取出递与士隐。}
}
{
    \eA{But on noticing the block newly-polished and brilliantly clear, which had
        moreover contracted in dimensions, and become no larger than the pendant of
        a fan, they were greatly filled with admiration. The Buddhist priest picked
        it up, and laid it in the palm of his hand.}
}
{
}

\Verse{16}
{
    \zA{士隐接了看时，原来是块鲜明美玉，上面字迹分明，镌着“通灵宝玉”四字，后面 还有几行小字。 正欲细看时，那僧便说“已到幻境”，就强从手中夺了去，和那道
        人竟过了一座大石牌坊，上面大书四字，乃是“太虚幻境”。 两边又有一副对联道： 假作真时真亦假， 无为有处有还无。
        士隐意欲也跟着过去，方举步时，忽听一声霹雳若山崩地陷，士隐大叫一声， 定睛看时，只见烈日炎炎，芭蕉冉冉，梦中之事便忘了一半。 又见奶母抱了英莲走 来。}
}
{
    \eA{"Your appearance," he said laughingly, "may well declare you to be a
        supernatural object, but as you lack any inherent quality it is necessary to
        inscribe a few characters on you, so that every one who shall see you may at
        once recognise you to be a remarkable thing.}
}
{
}

\Verse{17}
{
    \zA{士隐见女儿越发生得粉装玉琢，乖觉可喜，便伸手接来抱在怀中斗他玩耍一回； 又带至街前，看那过会的热闹。 方欲进来时，只见从那边来了一僧一道。 那僧癞头
        跣足，那道跛足蓬头，疯疯癫癫，挥霍谈笑而至。 及到了他门前，看见士隐抱着英 莲，那僧便大哭起来，又向士隐道：“施主，你把这有命无运、累及爹娘之物抱在
        怀内作甚！” 士隐听了，知是疯话，也不睬他。 那僧还说：“舍我罢!舍我罢！” 士隐 不耐烦，便抱着女儿转身。}
}
{
    \eA{And subsequently, when you will be taken into a country where honour and
        affluence will reign, into a family cultured in mind and of official status,
        in a land where flowers and trees shall flourish with luxuriance, in a town
        of refinement, renown and glory; when you once will have been there..." The
        stone listened with intense delight. "What characters may I ask," it
        consequently inquired, "will you inscribe?}
}
{
}

\Verse{18}
{
    \zA{才要进去，那僧乃指着他大笑，口内念了四句言词，道 是： 惯养娇生笑你痴，菱花空对雪澌澌。 好防佳节元宵后，便是烟消火灭时。
        士隐听得明白，心下犹豫，意欲问他来历。 只听道人说道：“你我不必同行，就此 分手，各干营生去罢。 三劫后我在北邙山等你，会齐了同往太虚幻境销号。” 那僧
        道：“最妙，最妙！” 说毕，二人一去，再不见个踪影了。 士隐心中此时自忖：这两个人必有来历，很该问他一问，如今后悔却已晚了。}
}
{
    \eA{and what place will I be taken to? pray, pray explain to me in lucid terms."
        "You mustn't be inquisitive," the bonze replied, with a smile, "in days to
        come you'll certainly understand everything." Having concluded these words,
        he forthwith put the stone in his sleeve, and proceeded leisurely on his
        journey, in company with the Taoist priest. Whither, however, he took the
        stone, is not divulged.}
}
{
}

\Verse{19}
{
    \zA{这士隐正在痴想，忽见隔壁葫芦庙内寄居的一个穷儒，姓贾名化、表字时飞、别号 雨村的走来。 这贾雨村原系湖州人氏，也是诗书仕宦之族。 因他生于末世，父母祖
        宗根基已尽，人口衰丧，只剩得他一身一口。 在家乡无益，因进京求取功名，再整 基业。 自前岁来此，又淹蹇住了，暂寄庙中安身，每日卖文作字为生，故士隐常与
        他交接。 当下雨村见了士隐，忙施礼陪笑道：“老先生倚门伫望，敢街市上有甚新 闻么？” 士隐笑道：“非也。}
}
{
    \eA{Nor can it be known how many centuries and ages elapsed, before a Taoist
        priest, K'ung K'ung by name, passed, during his researches after the eternal
        reason and his quest after immortality, by these Ta Huang Hills, Wu Ch'i
        cave and Ch'ing Keng Peak.}
}
{
}

\Verse{20}
{
    \zA{适因小女啼哭，引他出来作耍，正是无聊的很。 贾兄 来得正好，请入小斋，彼此俱可消此永昼。” 说着便令人送女儿进去，自携了雨村 来至书房中，小童献茶。
        方谈得三五句话，忽家人飞报：“严老爷来拜。” 士隐慌忙 起身谢道：“恕诓驾之罪，且请略坐，弟即来奉陪。” 雨村起身也让道：“老先生请 便。
        晚生乃常造之客，稍候何妨。” 说着士隐已出前厅去了。 这里雨村且翻弄诗籍解闷，忽听得窗外有女子嗽声。}
}
{
    \eA{Suddenly perceiving a large block of stone, on the surface of which the
        traces of characters giving, in a connected form, the various incidents of
        its fate, could be clearly deciphered, K'ung K'ung examined them from first
        to last.}
}
{
}

\Verse{21}
{
    \zA{雨村遂起身往外一看，原 来是一个丫鬟在那里掐花儿，生的仪容不俗，眉目清秀，虽无十分姿色，却也有动 人之处。 雨村不觉看得呆了。
        那甄家丫鬟掐了花儿方欲走时，猛抬头见窗内有人： 敝巾旧服，虽是贫窘，然生得腰圆背厚，面阔口方，更兼剑眉星眼，直鼻方腮。 这
        丫鬟忙转身回避，心下自想：“这人生的这样雄壮，却又这样褴褛，我家并无这样 贫窘亲友。 想他定是主人常说的什么贾雨村了，怪道又说他‘必非久困之人，每每
        有意帮助周济他，只是没什么机会。’”}
}
{
    \eA{They, in fact, explained how that this block of worthless stone had
        originally been devoid of the properties essential for the repairs to the
        heavens, how it would be transmuted into human form and introduced by Mang
        Mang the High Lord, and Miao Miao, the Divine, into the world of mortals,
        and how it would be led over the other bank (across the San Sara).}
}
{
}

\Verse{22}
{
    \zA{如此一想，不免又回头一两次。 雨村见他回 头，便以为这女子心中有意于他，遂狂喜不禁，自谓此女子必是个巨眼英豪、风尘 中之知己。
        一时小童进来，雨村打听得前面留饭，不可久待，遂从夹道中自便门出 去了。 士隐待客既散，知雨村已去，便也不去再邀。
        一日到了中秋佳节，士隐家宴已毕，又另具一席于书房，自己步月至庙中来邀 雨村。 原来雨村自那日见了甄家丫鬟曾回顾他两次，自谓是个知己，便时刻放在心 上。}
}
{
    \eA{On the surface, the record of the spot where it would fall, the place of its
        birth, as well as various family trifles and trivial love affairs of young
        ladies, verses, odes, speeches and enigmas was still complete; but the name
        of the dynasty and the year of the reign were obliterated, and could not be
        ascertained.}
}
{
}

\Verse{23}
{
    \zA{今又正值中秋，不免对月有怀，因而口占五言一律云： 未卜三生愿，频添一段愁。 闷来时敛额，行去几回眸。 自顾风前影，谁堪月下俦？ 蟾光如有意，先上玉人头。
        雨村吟罢，因又思及平生抱负，苦未逢时，乃又搔首对天长叹，复高吟一联云： 玉在椟中求善价，钗于奁内待时飞。
        恰值士隐走来听见，笑道：“雨村兄真抱负不凡也！” 雨村忙笑道：“不敢，不 过偶吟前人之句，何期过誉如此。” 因问：“老先生何兴至此？”}
}
{
    \eA{On the obverse, were also the following enigmatical verses:  Lacking in
        virtues meet the azure skies to mend,  In vain the mortal world full many a
        year I wend,  Of a former and after life these facts that be,  Who will for
        a tradition strange record for me? K'ung K'ung, the Taoist, having pondered
        over these lines for a while, became aware that this stone had a history of
        some kind.}
}
{
}

\Verse{24}
{
    \zA{士隐笑道：“今夜 中秋，俗谓团圆之节，想尊兄旅寄僧房，不无寂寥之感。 故特具小酌邀兄到敝斋一 饮，不知可纳芹意否？”
        雨村听了，并不推辞，便笑道：“既蒙谬爱，何敢拂此盛 情。” 说着便同士隐复过这边书院中来了。 须臾茶毕，早已设下杯盘，那美酒佳肴自不必说。
        二人归坐，先是款酌慢饮， 渐次谈至兴浓，不觉飞觥献？ 起来。 当时街坊上家家箫管，户户笙歌，当头一轮明 月，飞彩凝辉。 二人愈添豪兴，酒到杯干。}
}
{
    \eA{"Brother stone," he forthwith said, addressing the stone, "the concerns of
        past days recorded on you possess, according to your own account, a
        considerable amount of interest, and have been for this reason inscribed,
        with the intent of soliciting generations to hand them down as remarkable
        occurrences.}
}
{
}

\Verse{25}
{
    \zA{雨村此时已有七八分酒意，狂兴不禁， 乃对月寓怀，口占一绝云： 时逢三五便团？ ，满把清光护玉栏。 天上一轮才捧出，人间万姓仰头看。
        士隐听了大叫：“妙极!弟每谓兄必非久居人下者，今所吟之句，飞腾之兆已见，不 日可接履于云霄之上了。 可贺可贺！” 乃亲斟一斗为贺。
        雨村饮干，忽叹道：“非晚 生酒后狂言，若论时尚之学，晚生也或可去充数挂名。 只是如今行李路费一概无措， 神京路远，非赖卖字撰文即能到得。”}
}
{
    \eA{But in my own opinion, they lack, in the first place, any data by means of
        which to establish the name of the Emperor and the year of his reign; and,
        in the second place, these constitute no record of any excellent policy,
        adopted by any high worthies or high loyal statesmen, in the government of
        the state, or in the rule of public morals.}
}
{
}

\Verse{26}
{
    \zA{士隐不待说完，便道：“兄何不早言!弟已久 有此意，但每遇兄时并未谈及，故未敢唐突。 今既如此，弟虽不才：‘义利’二字 却还识得；
        且喜明岁正当大比，兄宜作速入都，春闱一捷，方不负兄之所学。 其盘 费馀事弟自代为处置，亦不枉兄之谬识矣。” 当下即命小童进去速封五十两白银并
        两套冬衣，又云：“十九日乃黄道之期，兄可即买舟西上。 待雄飞高举，明冬再晤， 岂非大快之事！” 雨村收了银衣，不过略谢一语，并不介意，仍是吃酒谈笑。}
}
{
    \eA{The contents simply treat of a certain number of maidens, of exceptional
        character; either of their love affairs or infatuations, or of their small
        deserts or insignificant talents; and were I to transcribe the whole
        collection of them, they would, nevertheless, not be estimated as a book of
        any exceptional worth." "Sir Priest," the stone replied with assurance, "why
        are you so excessively dull?}
}
{
}

\Verse{27}
{
    \zA{那天 已交三鼓，二人方散。 士隐送雨村去后，回房一觉，直至红日三竿方醒。 因思昨夜之事，意欲写荐书
        两封与雨村带至都中去，使雨村投谒个仕宦之家为寄身之地。 因使人过去请时，那 家人回来说：“和尚说，贾爷今日五鼓已进京去了，也曾留下话与和尚转达老爷，
        说：‘读书人不在黄道黑道，总以事理为要，不及面辞了。’” 士隐听了，也只得罢 了。 真是闲处光阴易过，倏忽又是元宵佳节。
        士隐令家人霍启抱了英莲，去看社火 花灯。}
}
{
    \eA{The dynasties recorded in the rustic histories, which have been written from
        age to age, have, I am fain to think, invariably assumed, under false
        pretences, the mere nomenclature of the Han and T'ang dynasties.}
}
{
}

\Verse{28}
{
    \zA{半夜中霍启因要小解，便将英莲放在一家门槛上坐着。 待他小解完了来抱时， 那有英莲的踪影？ 急的霍启直寻了半夜。 至天明不见，那霍启也不敢回来见主人，
        便逃往他乡去了。 那士隐夫妇见女儿一夜不归，便知有些不好； 再使几人去找寻， 回来皆云影响全无。
        夫妻二人半世只生此女，一旦失去，何等烦恼，因此昼夜啼哭， 几乎不顾性命。 看看一月，士隐已先得病，夫人封氏也因思女构疾，日日请医问卦。}
}
{
    \eA{They differ from the events inscribed on my block, which do not borrow this
        customary practice, but, being based on my own experiences and natural
        feelings, present, on the contrary, a novel and unique character.}
}
{
}

\Verse{29}
{
    \zA{不想这日 三月十五，葫芦庙中炸供，那和尚不小心，油锅火逸，便烧着窗纸。 此方人家俱用
        竹篱木壁，也是劫数应当如此，于是接二连三牵五挂四，将一条街烧得如火焰山一 般。 彼时虽有军民来救，那火已成了势了，如何救得下？ 直烧了一夜方息，也不知
        烧了多少人家。 只可怜甄家在隔壁，早成了一堆瓦砾场了，只有他夫妇并几个家人 的性命不曾伤了。 急的士隐惟跌足长叹而已。 与妻子商议，且到田庄上去住。}
}
{
    \eA{Besides, in the pages of these rustic histories, either the aspersions upon
        sovereigns and statesmen, or the strictures upon individuals, their wives,
        and their daughters, or the deeds of licentiousness and violence are too
        numerous to be computed. Indeed, there is one more kind of loose literature,
        the wantonness and pollution in which work most easy havoc upon youth.}
}
{
}

\Verse{30}
{
    \zA{偏值 近年水旱不收，贼盗蜂起，官兵剿捕，田庄上又难以安身，只得将田地都折变了， 携了妻子与两个丫鬟投他岳丈家去。
        他岳丈名唤封肃，本贯大如州人氏，虽是务农，家中却还殷实。 今见女婿这等 狼狈而来，心中便有些不乐。 幸而士隐还有折变田产的银子在身边，拿出来托他随
        便置买些房地，以为后日衣食之计，那封肃便半用半赚的，略与他些薄田破屋。 士 隐乃读书之人，不惯生理稼穑等事，勉强支持了一二年，越发穷了。 封肃见面时，
        便说些现成话儿；}
}
{
    \eA{"As regards the works, in which the characters of scholars and beauties is
        delineated their allusions are again repeatedly of Wen Chün, their theme in
        every page of Tzu Chien; a thousand volumes present no diversity; and a
        thousand characters are but a counterpart of each other. What is more, these
        works, throughout all their pages, cannot help bordering on extreme licence.}
}
{
}

\Verse{31}
{
    \zA{且人前人后又怨他不会过，只一味好吃懒做。 士隐知道了，心中 未免悔恨，再兼上年惊唬，急忿怨痛，暮年之人，那禁得贫病交攻，竟渐渐的露出 了那下世的光景来。
        可巧这日拄了拐扎挣到街前散散心时，忽见那边来了一个跛足道人，疯狂落拓， 麻鞋鹑衣，口内念着几句言词道： 世人都晓神仙好，惟有功名忘不了。
        古今将相在何方？ 荒冢一堆草没了。 世人都晓神仙好，只有金银忘不了。 终朝只恨聚无多，及到多时眼闭了。 世人都晓神仙好，只有娇妻忘不了。}
}
{
    \eA{The authors, however, had no other object in view than to give utterance to
        a few sentimental odes and elegant ballads of their own, and for this reason
        they have fictitiously invented the names and surnames of both men and
        women, and necessarily introduced, in addition, some low characters, who
        should, like a buffoon in a play, create some excitement in the plot.}
}
{
}

\Verse{32}
{
    \zA{君生日日说恩情，君死又随人去了。 世人都晓神仙好，只有儿孙忘不了。 痴心父母古来多，孝顺子孙谁见了？ 士隐听了，便迎上来道：“你满口说些什么？
        只听见些‘好’‘了’‘好’‘了’。” 那 道人笑道：“你若果听见‘好’‘了’二字，还算你明白：可知世上万般，好便是了， 了便是好。 若不了，便不好；
        若要好，须是了。 我这歌儿便叫《好了歌》。” 士隐本 是有夙慧的，一闻此言，心中早已悟彻，因笑道：“且住，待我将你这《好了歌》 注解出来何如？”}
}
{
    \eA{"Still more loathsome is a kind of pedantic and profligate literature,
        perfectly devoid of all natural sentiment, full of self-contradictions; and,
        in fact, the contrast to those maidens in my work, whom I have, during half
        my lifetime, seen with my own eyes and heard with my own ears.}
}
{
}

\Verse{33}
{
    \zA{道人笑道：“你就请解。” 士隐乃说道： 陋室空堂，当年笏满床。 衰草枯杨，曾为歌舞场。 蛛丝儿结满雕粱，绿纱今又 在蓬窗上。
        说甚么脂正浓、粉正香，如何两鬓又成霜？ 昨日黄土陇头埋白骨，今宵 红绡帐底卧鸳鸯。 金满箱，银满箱，转眼乞丐人皆谤。 正叹他人命不长，那知自己归来丧？
        训有方， 保不定日后作强梁。 择膏粱，谁承望流落在烟花巷!因嫌纱帽小，致使锁枷扛。 昨 怜破袄寒，今嫌紫蟒长：乱烘烘你方唱罢我登场，反认他乡是故乡。}
}
{
    \eA{And though I will not presume to estimate them as superior to the heroes and
        heroines in the works of former ages, yet the perusal of the motives and
        issues of their experiences, may likewise afford matter sufficient to banish
        dulness, and to break the spell of melancholy.}
}
{
}

\Verse{34}
{
    \zA{甚荒唐，到头 来都是“为他人作嫁衣裳”。 那疯跛道人听了，拍掌大笑道：“解得切!解得切！” 士隐便说一声“走罢”，将
        道人肩上的搭裢抢过来背上，竟不回家，同着疯道人飘飘而去。 当下哄动街坊，众 人当作一件新闻传说。 封氏闻知此信，哭个死去活来。 只得与父亲商议，遣人各处
        访寻，那讨音信？ 无奈何，只得依靠着他父母度日。 幸而身边还有两个旧日的丫鬟 伏侍，主仆三人，日夜作些针线，帮着父亲用度。 那封肃虽然每日抱怨，也无可奈
        何了。}
}
{
    \eA{"As regards the several stanzas of doggerel verse, they may too evoke such
        laughter as to compel the reader to blurt out the rice, and to spurt out the
        wine.}
}
{
}

\Verse{35}
{
    \zA{这日那甄家的大丫鬟在门前买线，忽听得街上喝道之声。 众人都说：“新太爷 到任了！” 丫鬟隐在门内看时，只见军牢快手一对一对过去，俄而大轿内抬着一个
        乌帽猩袍的官府来了。 那丫鬟倒发了个怔，自思：“这官儿好面善？ 倒像在那里见过 的。” 于是进入房中，也就丢过不在心上。
        至晚间正待歇息之时，忽听一片声打的 门响，许多人乱嚷，说：“本县太爷的差人来传人问话！” 封肃听了，唬得目瞪口呆。 不知有何祸事，且听下回分解。
        (本章完)}
}
{
    \eA{"In these pages, the scenes depicting the anguish of separation, the bliss
        of reunion, and the fortunes of prosperity and of adversity are all, in
        every detail, true to human nature, and I have not taken upon myself to make
        the slightest addition, or alteration, which might lead to the perversion of
        the truth.}
}
{
}

\Verse{36}
{
    \zA{}
}
{
    \eA{"My only object has been that men may, after a drinking bout, or after they
        wake from sleep or when in need of relaxation from the pressure of business,
        take up this light literature, and not only expunge the traces of antiquated
        books, and obtain a new kind of distraction, but that they may also lay by a
        long life as well as energy and strength; for it bears no point of
        similarity to those works, whose designs are false, whose course is immoral.
        Now, Sir Priest, what are your views on the subject?" K'ung K'ung having
        pondered for a while over the words, to which he had listened intently,
        re-perused, throughout, this record of the stone;}
}
{
}

\Verse{37}
{
    \zA{}
}
{
    \eA{and finding that the general purport consisted of nought else than a
        treatise on love, and likewise of an accurate transcription of facts,
        without the least taint of profligacy injurious to the times, he thereupon
        copied the contents, from beginning to end, to the intent of charging the
        world to hand them down as a strange story. Hence it was that K'ung K'ung,
        the Taoist, in consequence of his perception, (in his state of) abstraction,
        of passion, the generation, from this passion, of voluptuousness, the
        transmission of this voluptuousness into passion, and the apprehension, by
        means of passion, of its unreality, forthwith altered his name for that of
        "Ch'ing Tseng" (the Voluptuous Bonze), and changed the title of "the Memoir
        of a Stone" (Shih-t'ou-chi,) for that of "Ch'ing Tseng Lu," The Record of
        the Voluptuous Bonze;}
}
{
}

\Verse{38}
{
    \zA{}
}
{
    \eA{while K'ung Mei-chi of Tung Lu gave it the name of "Feng Yüeh Pao Chien,"
        "The Precious Mirror of Voluptuousness." In later years, owing to the
        devotion by Tsao Hsüeh-ch'in in the Tao Hung study, of ten years to the
        perusal and revision of the work, the additions and modifications effected
        by him five times, the affix of an index and the division into periods and
        chapters, the book was again entitled "Chin Ling Shih Erh Ch'ai," "The
        Twelve Maidens of Chin Ling." A stanza was furthermore composed for the
        purpose.}
}
{
}

\Verse{39}
{
    \zA{}
}
{
    \eA{This then, and no other, is the origin of the Record of the Stone. The poet
        says appositely:--  Pages full of silly litter,  Tears a handful sour and
        bitter; All a fool the author hold,  But their zest who can unfold? You have
        now understood the causes which brought about the Record of the Stone, but
        as you are not, as yet, aware what characters are depicted, and what
        circumstances are related on the surface of the block, reader, please lend
        an ear to the narrative on the stone, which runs as follows:-- In old days,
        the land in the South East lay low. In this South-East part of the world,
        was situated a walled town, Ku Su by name.}
}
{
}

\Verse{40}
{
    \zA{}
}
{
    \eA{Within the walls a locality, called the Ch'ang Men, was more than all others
        throughout the mortal world, the centre, which held the second, if not the
        first place for fashion and life. Beyond this Ch'ang Men was a street called
        Shih-li-chieh (Ten \_Li\_ street); in this street a lane, the Jen Ch'ing
        lane (Humanity and Purity); and in this lane stood an old temple, which on
        account of its diminutive dimensions, was called, by general consent, the
        Gourd temple. Next door to this temple lived the family of a district
        official, Chen by surname, Fei by name, and Shih-yin by style. His wife, née
        Feng, possessed a worthy and virtuous disposition, and had a clear
        perception of moral propriety and good conduct.}
}
{
}

\Verse{41}
{
    \zA{}
}
{
    \eA{This family, though not in actual possession of excessive affluence and
        honours, was, nevertheless, in their district, conceded to be a clan of
        well-to-do standing. As this Chen Shih-yin was of a contented and
        unambitious frame of mind, and entertained no hankering after any official
        distinction, but day after day of his life took delight in gazing at
        flowers, planting bamboos, sipping his wine and conning poetical works, he
        was in fact, in the indulgence of these pursuits, as happy as a supernatural
        being. One thing alone marred his happiness.}
}
{
}

\Verse{42}
{
    \zA{}
}
{
    \eA{He had lived over half a century and had, as yet, no male offspring around
        his knees. He had one only child, a daughter, whose infant name was Ying
        Lien. She was just three years of age. On a long summer day, on which the
        heat had been intense, Shih-yin sat leisurely in his library. Feeling his
        hand tired, he dropped the book he held, leant his head on a teapoy, and
        fell asleep. Of a sudden, while in this state of unconsciousness, it seemed
        as if he had betaken himself on foot to some spot or other whither he could
        not discriminate. Unexpectedly he espied, in the opposite direction, two
        priests coming towards him: the one a Buddhist, the other a Taoist. As they
        advanced they kept up the conversation in which they were engaged. "Whither
        do you purpose taking the object you have brought away?"}
}
{
}

\Verse{43}
{
    \zA{}
}
{
    \eA{he heard the Taoist inquire. To this question the Buddhist replied with a
        smile: "Set your mind at ease," he said; "there's now in maturity a plot of
        a general character involving mundane pleasures, which will presently come
        to a denouement. The whole number of the votaries of voluptuousness have, as
        yet, not been quickened or entered the world, and I mean to avail myself of
        this occasion to introduce this object among their number, so as to give it
        a chance to go through the span of human existence."}
}
{
}

\Verse{44}
{
    \zA{}
}
{
    \eA{"The votaries of voluptuousness of these days will naturally have again to
        endure the ills of life during their course through the mortal world," the
        Taoist remarked; "but when, I wonder, will they spring into existence? and
        in what place will they descend?" "The account of these circumstances," the
        bonze ventured to reply, "is enough to make you laugh! They amount to this:
        there existed in the west, on the bank of the Ling (spiritual) river, by the
        side of the San Sheng (thrice-born) stone, a blade of the Chiang Chu (purple
        pearl) grass. At about the same time it was that the block of stone was,
        consequent upon its rejection by the goddess of works, also left to ramble
        and wander to its own gratification, and to roam about at pleasure to every
        and any place.}
}
{
}

\Verse{45}
{
    \zA{}
}
{
    \eA{One day it came within the precincts of the Ching Huan (Monitory Vision)
        Fairy; and this Fairy, cognizant of the fact that this stone had a history,
        detained it, therefore, to reside at the Ch'ih Hsia (purple clouds) palace,
        and apportioned to it the duties of attendant on Shen Ying, a fairy of the
        Ch'ih Hsia palace. "This stone would, however, often stroll along the banks
        of the Ling river, and having at the sight of the blade of spiritual grass
        been filled with admiration, it, day by day, moistened its roots with sweet
        dew. This purple pearl grass, at the outset, tarried for months and years;}
}
{
}

\Verse{46}
{
    \zA{}
}
{
    \eA{but being at a later period imbued with the essence and luxuriance of heaven
        and earth, and having incessantly received the moisture and nurture of the
        sweet dew, divested itself, in course of time, of the form of a grass;
        assuming, in lieu, a human nature, which gradually became perfected into the
        person of a girl. "Every day she was wont to wander beyond the confines of
        the Li Hen (divested animosities) heavens. When hungry she fed on the Pi
        Ch'ing (hidden love) fruit--when thirsty she drank the Kuan ch'ou
        (discharged sorrows,) water. Having, however, up to this time, not shewn her
        gratitude for the virtue of nurture lavished upon her, the result was but
        natural that she should resolve in her heart upon a constant and incessant
        purpose to make suitable acknowledgment.}
}
{
}

\Verse{47}
{
    \zA{}
}
{
    \eA{"I have been," she would often commune within herself, "the recipient of the
        gracious bounty of rain and dew, but I possess no such water as was lavished
        upon me to repay it! But should it ever descend into the world in the form
        of a human being, I will also betake myself thither, along with it; and if I
        can only have the means of making restitution to it, with the tears of a
        whole lifetime, I may be able to make adequate return." "This resolution it
        is that will evolve the descent into the world of so many pleasure-bound
        spirits of retribution and the experience of fantastic destinies; and this
        crimson pearl blade will also be among the number.}
}
{
}

\Verse{48}
{
    \zA{}
}
{
    \eA{The stone still lies in its original place, and why should not you and I
        take it along before the tribunal of the Monitory Vision Fairy, and place on
        its behalf its name on record, so that it should descend into the world, in
        company with these spirits of passion, and bring this plot to an issue?" "It
        is indeed ridiculous," interposed the Taoist. "Never before have I heard
        even the very mention of restitution by means of tears! Why should not you
        and I avail ourselves of this opportunity to likewise go down into the
        world? and if successful in effecting the salvation of a few of them, will
        it not be a work meritorious and virtuous?" "This proposal," remarked the
        Buddhist, "is quite in harmony with my own views.}
}
{
}

\Verse{49}
{
    \zA{}
}
{
    \eA{Come along then with me to the palace of the Monitory Vision Fairy, and let
        us deliver up this good-for-nothing object, and have done with it! And when
        the company of pleasure-bound spirits of wrath descend into human existence,
        you and I can then enter the world. Half of them have already fallen into
        the dusty universe, but the whole number of them have not, as yet, come
        together." "Such being the case," the Taoist acquiesced, "I am ready to
        follow you, whenever you please to go." But to return to Chen Shih-yin.
        Having heard every one of these words distinctly, he could not refrain from
        forthwith stepping forward and paying homage. "My spiritual lords," he said,
        as he smiled, "accept my obeisance."}
}
{
}

\Verse{50}
{
    \zA{}
}
{
    \eA{The Buddhist and Taoist priests lost no time in responding to the
        compliment, and they exchanged the usual salutations. "My spiritual lords,"
        Shih-yin continued; "I have just heard the conversation that passed between
        you, on causes and effects, a conversation the like of which few mortals
        have forsooth listened to; but your younger brother is sluggish of
        intellect, and cannot lucidly fathom the import! Yet could this dulness and
        simplicity be graciously dispelled, your younger brother may, by listening
        minutely, with undefiled ear and careful attention, to a certain degree be
        aroused to a sense of understanding; and what is more, possibly find the
        means of escaping the anguish of sinking down into Hades."}
}
{
}

\Verse{51}
{
    \zA{}
}
{
    \eA{The two spirits smiled, "The conversation," they added, "refers to the
        primordial scheme and cannot be divulged before the proper season; but, when
        the time comes, mind do not forget us two, and you will readily be able to
        escape from the fiery furnace." Shih-yin, after this reply, felt it
        difficult to make any further inquiries. "The primordial scheme," he however
        remarked smiling, "cannot, of course, be divulged; but what manner of thing,
        I wonder, is the good-for-nothing object you alluded to a short while back?
        May I not be allowed to judge for myself?" "This object about which you
        ask," the Buddhist Bonze responded, "is intended, I may tell you, by fate to
        be just glanced at by you." With these words he produced it, and handed it
        over to Shih-yin. Shih-yin received it.}
}
{
}

\Verse{52}
{
    \zA{}
}
{
    \eA{On scrutiny he found it, in fact, to be a beautiful gem, so lustrous and so
        clear that the traces of characters on the surface were distinctly visible.
        The characters inscribed consisted of the four "T'ung Ling Pao Yü,"
        "Precious Gem of Spiritual Perception." On the obverse, were also several
        columns of minute words, which he was just in the act of looking at
        intently, when the Buddhist at once expostulated. "We have already reached,"
        he exclaimed, "the confines of vision." Snatching it violently out of his
        hands, he walked away with the Taoist, under a lofty stone portal, on the
        face of which appeared in large type the four characters: "T'ai Hsü Huan
        Ching," "The Visionary limits of the Great Void."}
}
{
}

\Verse{53}
{
    \zA{}
}
{
    \eA{On each side was a scroll with the lines:  When falsehood stands for truth,
        truth likewise becomes false,  Where naught be made to aught, aught changes
        into naught. Shih-yin meant also to follow them on the other side, but, as
        he was about to make one step forward, he suddenly heard a crash, just as if
        the mountains had fallen into ruins, and the earth sunk into destruction. As
        Shih-yin uttered a loud shout, he looked with strained eye; but all he could
        see was the fiery sun shining, with glowing rays, while the banana leaves
        drooped their heads. By that time, half of the circumstances connected with
        the dream he had had, had already slipped from his memory. He also noticed a
        nurse coming towards him with Ying Lien in her arms.}
}
{
}

\Verse{54}
{
    \zA{}
}
{
    \eA{To Shih-yin's eyes his daughter appeared even more beautiful, such a bright
        gem, so precious, and so lovable. Forthwith stretching out his arms, he took
        her over, and, as he held her in his embrace, he coaxed her to play with him
        for a while; after which he brought her up to the street to see the great
        stir occasioned by the procession that was going past. He was about to come
        in, when he caught sight of two priests, one a Taoist, the other a Buddhist,
        coming hither from the opposite direction. The Buddhist had a head covered
        with mange, and went barefooted. The Taoist had a limping foot, and his hair
        was all dishevelled. Like maniacs, they jostled along, chattering and
        laughing as they drew near.}
}
{
}

\Verse{55}
{
    \zA{}
}
{
    \eA{As soon as they reached Shih-yin's door, and they perceived him with Ying
        Lien in his arms, the Bonze began to weep aloud. Turning towards Shih-yin,
        he said to him: "My good Sir, why need you carry in your embrace this living
        but luckless thing, which will involve father and mother in trouble?" These
        words did not escape Shih-yin's ear; but persuaded that they amounted to
        raving talk, he paid no heed whatever to the bonze. "Part with her and give
        her to me," the Buddhist still went on to say. Shih-yin could not restrain
        his annoyance; and hastily pressing his daughter closer to him, he was
        intent upon going in, when the bonze pointed his hand at him, and burst out
        in a loud fit of laughter.}
}
{
}

\Verse{56}
{
    \zA{}
}
{
    \eA{He then gave utterance to the four lines that follow:  You indulge your
        tender daughter and are laughed at as inane; Vain you face the snow, oh
        mirror! for it will evanescent wane,  When the festival of lanterns is gone
        by, guard 'gainst your doom,  'Tis what time the flames will kindle, and the
        fire will consume. Shih-yin understood distinctly the full import of what he
        heard; but his heart was still full of conjectures. He was about to inquire
        who and what they were, when he heard the Taoist remark,--"You and I cannot
        speed together; let us now part company, and each of us will be then able to
        go after his own business. After the lapse of three ages, I shall be at the
        Pei Mang mount, waiting for you;}
}
{
}

\Verse{57}
{
    \zA{}
}
{
    \eA{and we can, after our reunion, betake ourselves to the Visionary Confines of
        the Great Void, there to cancel the name of the stone from the records."
        "Excellent! first rate!" exclaimed the Bonze. And at the conclusion of these
        words, the two men parted, each going his own way, and no trace was again
        seen of them. "These two men," Shih-yin then pondered within his heart,
        "must have had many experiences, and I ought really to have made more
        inquiries of them; but at this juncture to indulge in regret is anyhow too
        late."}
}
{
}

\Verse{58}
{
    \zA{}
}
{
    \eA{While Shih-yin gave way to these foolish reflections, he suddenly noticed
        the arrival of a penniless scholar, Chia by surname, Hua by name, Shih-fei
        by style and Yü-ts'un by nickname, who had taken up his quarters in the
        Gourd temple next door. This Chia Yü-ts'un was originally a denizen of
        Hu-Chow, and was also of literary and official parentage, but as he was born
        of the youngest stock, and the possessions of his paternal and maternal
        ancestors were completely exhausted, and his parents and relatives were
        dead, he remained the sole and only survivor;}
}
{
}

\Verse{59}
{
    \zA{}
}
{
    \eA{and, as he found his residence in his native place of no avail, he therefore
        entered the capital in search of that reputation, which would enable him to
        put the family estate on a proper standing. He had arrived at this place
        since the year before last, and had, what is more, lived all along in very
        straitened circumstances. He had made the temple his temporary quarters, and
        earned a living by daily occupying himself in composing documents and
        writing letters for customers. Thus it was that Shih-yin had been in
        constant relations with him. As soon as Yü-ts'un perceived Shih-yin, he lost
        no time in saluting him. "My worthy Sir," he observed with a forced smile;
        "how is it you are leaning against the door and looking out?}
}
{
}

\Verse{60}
{
    \zA{}
}
{
    \eA{Is there perchance any news astir in the streets, or in the public places?"
        "None whatever," replied Shih-yin, as he returned the smile. "Just a while
        back, my young daughter was in sobs, and I coaxed her out here to amuse her.
        I am just now without anything whatever to attend to, so that, dear brother
        Chia, you come just in the nick of time. Please walk into my mean abode, and
        let us endeavour, in each other's company, to while away this long summer
        day." After he had made this remark, he bade a servant take his daughter in,
        while he, hand-in-hand with Yü-ts'un, walked into the library, where a young
        page served tea.}
}
{
}

\Verse{61}
{
    \zA{}
}
{
    \eA{They had hardly exchanged a few sentences, when one of the household came
        in, in flying haste, to announce that Mr. Yen had come to pay a visit.
        Shih-yin at once stood up. "Pray excuse my rudeness," he remarked
        apologetically, "but do sit down; I shall shortly rejoin you, and enjoy the
        pleasure of your society." "My dear Sir," answered Yü-ts'un, as he got up,
        also in a conceding way, "suit your own convenience. I've often had the
        honour of being your guest, and what will it matter if I wait a little?"
        While these apologies were yet being spoken, Shih-yin had already walked out
        into the front parlour. During his absence, Yü-ts'un occupied himself in
        turning over the pages of some poetical work to dispel ennui, when suddenly
        he heard, outside the window, a woman's cough. Yü-ts'un hurriedly got up and
        looked out.}
}
{
}

\Verse{62}
{
    \zA{}
}
{
    \eA{He saw at a glance that it was a servant girl engaged in picking flowers.
        Her deportment was out of the common; her eyes so bright, her eyebrows so
        well defined. Though not a perfect beauty, she possessed nevertheless charms
        sufficient to arouse the feelings. Yü-ts'un unwittingly gazed at her with
        fixed eye. This waiting-maid, belonging to the Chen family, had done picking
        flowers, and was on the point of going in, when she of a sudden raised her
        eyes and became aware of the presence of some person inside the window,
        whose head-gear consisted of a turban in tatters, while his clothes were the
        worse for wear.}
}
{
}

\Verse{63}
{
    \zA{}
}
{
    \eA{But in spite of his poverty, he was naturally endowed with a round waist, a
        broad back, a fat face, a square mouth; added to this, his eyebrows were
        swordlike, his eyes resembled stars, his nose was straight, his cheeks
        square. This servant girl turned away in a hurry and made her escape. "This
        man so burly and strong," she communed within herself, "yet at the same time
        got up in such poor attire, must, I expect, be no one else than the man,
        whose name is Chia Yü-ts'un or such like, time after time referred to by my
        master, and to whom he has repeatedly wished to give a helping hand, but has
        failed to find a favourable opportunity.}
}
{
}

\Verse{64}
{
    \zA{}
}
{
    \eA{And as related to our family there is no connexion or friend in such
        straits, I feel certain it cannot be any other person than he. Strange to
        say, my master has further remarked that this man will, for a certainty, not
        always continue in such a state of destitution." As she indulged in this
        train of thought, she could not restrain herself from turning her head round
        once or twice. When Yü-ts'un perceived that she had looked back, he readily
        interpreted it as a sign that in her heart her thoughts had been of him, and
        he was frantic with irrepressible joy. "This girl," he mused, "is, no doubt,
        keen-eyed and eminently shrewd, and one in this world who has seen through
        me." The servant youth, after a short time, came into the room;}
}
{
}

\Verse{65}
{
    \zA{}
}
{
    \eA{and when Yü-ts'un made inquiries and found out from him that the guests in
        the front parlour had been detained to dinner, he could not very well wait
        any longer, and promptly walked away down a side passage and out of a back
        door. When the guests had taken their leave, Shih-yin did not go back to
        rejoin Yü-ts'un, as he had come to know that he had already left. In time
        the mid-autumn festivities drew near; and Shih-yin, after the family banquet
        was over, had a separate table laid in the library, and crossed over, in the
        moonlight, as far as the temple and invited Yü-ts'un to come round.}
}
{
}

\Verse{66}
{
    \zA{}
}
{
    \eA{The fact is that Yü-ts'un, ever since the day on which he had seen the girl
        of the Chen family turn twice round to glance at him, flattered himself that
        she was friendly disposed towards him, and incessantly fostered fond
        thoughts of her in his heart. And on this day, which happened to be the
        mid-autumn feast, he could not, as he gazed at the moon, refrain from
        cherishing her remembrance. Hence it was that he gave vent to these
        pentameter verses:  Alas! not yet divined my lifelong wish,  And anguish
        ceaseless comes upon anguish  I came, and sad at heart, my brow I frowned;
        She went, and oft her head to look turned round. Facing the breeze, her
        shadow she doth watch,  Who's meet this moonlight night with her to match?}
}
{
}

\Verse{67}
{
    \zA{}
}
{
    \eA{The lustrous rays if they my wish but read  Would soon alight upon her
        beauteous head! Yü-ts'un having, after this recitation, recalled again to
        mind how that throughout his lifetime his literary attainments had had an
        adverse fate and not met with an opportunity (of reaping distinction), went
        on to rub his brow, and as he raised his eyes to the skies, he heaved a deep
        sigh and once more intoned a couplet aloud:  The gem in the cask a high
        price it seeks,  The pin in the case to take wing it waits.}
}
{
}

\Verse{68}
{
    \zA{}
}
{
    \eA{As luck would have it, Shih-yin was at the moment approaching, and upon
        hearing the lines, he said with a smile: "My dear Yü-ts'un, really your
        attainments are of no ordinary capacity." Yü-ts'un lost no time in smiling
        and replying. "It would be presumption in my part to think so," he observed.
        "I was simply at random humming a few verses composed by former writers, and
        what reason is there to laud me to such an excessive degree? To what, my
        dear Sir, do I owe the pleasure of your visit?" he went on to inquire.
        "Tonight," replied Shih-yin, "is the mid-autumn feast, generally known as
        the full-moon festival;}
}
{
}

\Verse{69}
{
    \zA{}
}
{
    \eA{and as I could not help thinking that living, as you my worthy brother are,
        as a mere stranger in this Buddhist temple, you could not but experience the
        feeling of loneliness. I have, for the express purpose, prepared a small
        entertainment, and will be pleased if you will come to my mean abode to have
        a glass of wine. But I wonder whether you will entertain favourably my
        modest invitation?" Yü-ts'un, after listening to the proposal, put forward
        no refusal of any sort; but remarked complacently: "Being the recipient of
        such marked attention, how can I presume to repel your generous
        consideration?" As he gave expression to these words, he walked off there
        and then, in company with Shih-yin, and came over once again into the court
        in front of the library. In a few minutes, tea was over.}
}
{
}

\Verse{70}
{
    \zA{}
}
{
    \eA{The cups and dishes had been laid from an early hour, and needless to say
        the wines were luscious; the fare sumptuous. The two friends took their
        seats. At first they leisurely replenished their glasses, and quietly sipped
        their wine; but as, little by little, they entered into conversation, their
        good cheer grew more genial, and unawares the glasses began to fly round,
        and the cups to be exchanged. At this very hour, in every house of the
        neighbourhood, sounded the fife and lute, while the inmates indulged in
        music and singing. Above head, the orb of the radiant moon shone with an
        all-pervading splendour, and with a steady lustrous light, while the two
        friends, as their exuberance increased, drained their cups dry so soon as
        they reached their lips.}
}
{
}

\Verse{71}
{
    \zA{}
}
{
    \eA{Yü-ts'un, at this stage of the collation, was considerably under the
        influence of wine, and the vehemence of his high spirits was irrepressible.
        As he gazed at the moon, he fostered thoughts, to which he gave vent by the
        recital of a double couplet. 'Tis what time three meets five, Selene is a
        globe! Her pure rays fill the court, the jadelike rails enrobe! Lo! in the
        heavens her disk to view doth now arise,  And in the earth below to gaze men
        lift their eyes. "Excellent!" cried Shih-yin with a loud voice, after he had
        heard these lines; "I have repeatedly maintained that it was impossible for
        you to remain long inferior to any, and now the verses you have recited are
        a prognostic of your rapid advancement.}
}
{
}

\Verse{72}
{
    \zA{}
}
{
    \eA{Already it is evident that, before long, you will extend your footsteps far
        above the clouds! I must congratulate you! I must congratulate you! Let me,
        with my own hands, pour a glass of wine to pay you my compliments." Yü-ts'un
        drained the cup. "What I am about to say," he explained as he suddenly
        heaved a sigh, "is not the maudlin talk of a man under the effects of wine.
        As far as the subjects at present set in the examinations go, I could,
        perchance, also have well been able to enter the list, and to send in my
        name as a candidate; but I have, just now, no means whatever to make
        provision for luggage and for travelling expenses.}
}
{
}

\Verse{73}
{
    \zA{}
}
{
    \eA{The distance too to Shen Ching is a long one, and I could not depend upon
        the sale of papers or the composition of essays to find the means of getting
        there." Shih-yin gave him no time to conclude. "Why did you not speak about
        this sooner?" he interposed with haste. "I have long entertained this
        suspicion; but as, whenever I met you, this conversation was never broached,
        I did not presume to make myself officious. But if such be the state of
        affairs just now, I lack, I admit, literary qualification, but on the two
        subjects of friendly spirit and pecuniary means, I have, nevertheless, some
        experience. Moreover, I rejoice that next year is just the season for the
        triennial examinations, and you should start for the capital with all
        despatch;}
}
{
}

\Verse{74}
{
    \zA{}
}
{
    \eA{and in the tripos next spring, you will, by carrying the prize, be able to
        do justice to the proficiency you can boast of. As regards the travelling
        expenses and the other items, the provision of everything necessary for you
        by my own self will again not render nugatory your mean acquaintance with
        me." Forthwith, he directed a servant lad to go and pack up at once fifty
        taels of pure silver and two suits of winter clothes. "The nineteenth," he
        continued, "is a propitious day, and you should lose no time in hiring a
        boat and starting on your journey westwards.}
}
{
}

\Verse{75}
{
    \zA{}
}
{
    \eA{And when, by your eminent talents, you shall have soared high to a lofty
        position, and we meet again next winter, will not the occasion be extremely
        felicitous?" Yü-ts'un accepted the money and clothes with but scanty
        expression of gratitude. In fact, he paid no thought whatever to the gifts,
        but went on, again drinking his wine, as he chattered and laughed. It was
        only when the third watch of that day had already struck that the two
        friends parted company; and Shih-yin, after seeing Yü-ts'un off, retired to
        his room and slept, with one sleep all through, never waking until the sun
        was well up in the skies.}
}
{
}

\Verse{76}
{
    \zA{}
}
{
    \eA{Remembering the occurrence of the previous night, he meant to write a couple
        of letters of recommendation for Yü-ts'un to take along with him to the
        capital, to enable him, after handing them over at the mansions of certain
        officials, to find some place as a temporary home. He accordingly despatched
        a servant to ask him to come round, but the man returned and reported that
        from what the bonze said, "Mr. Chia had started on his journey to the
        capital, at the fifth watch of that very morning, that he had also left a
        message with the bonze to deliver to you, Sir, to the effect that men of
        letters paid no heed to lucky or unlucky days, that the sole consideration
        with them was the nature of the matter in hand, and that he could find no
        time to come round in person and bid good-bye."}
}
{
}

\Verse{77}
{
    \zA{}
}
{
    \eA{Shih-yin after hearing this message had no alternative but to banish the
        subject from his thoughts. In comfortable circumstances, time indeed goes by
        with easy stride. Soon drew near also the happy festival of the 15th of the
        1st moon, and Shih-yin told a servant Huo Ch'i to take Ying Lien to see the
        sacrificial fires and flowery lanterns. About the middle of the night, Huo
        Ch'i was hard pressed, and he forthwith set Ying Lien down on the doorstep
        of a certain house. When he felt relieved, he came back to take her up, but
        failed to find anywhere any trace of Ying Lien. In a terrible plight, Huo
        Ch'i prosecuted his search throughout half the night; but even by the dawn
        of day, he had not discovered any clue of her whereabouts.}
}
{
}

\Verse{78}
{
    \zA{}
}
{
    \eA{Huo Ch'i, lacking, on the other hand, the courage to go back and face his
        master, promptly made his escape to his native village. Shih-yin--in fact,
        the husband as well as the wife--seeing that their child had not come home
        during the whole night, readily concluded that some mishap must have
        befallen her. Hastily they despatched several servants to go in search of
        her, but one and all returned to report that there was neither vestige nor
        tidings of her. This couple had only had this child, and this at the
        meridian of their life, so that her sudden disappearance plunged them in
        such great distress that day and night they mourned her loss to such a point
        as to well nigh pay no heed to their very lives.}
}
{
}

\Verse{79}
{
    \zA{}
}
{
    \eA{A month in no time went by. Shih-yin was the first to fall ill, and his
        wife, Dame Feng, likewise, by dint of fretting for her daughter, was also
        prostrated with sickness. The doctor was, day after day, sent for, and the
        oracle consulted by means of divination. Little did any one think that on
        this day, being the 15th of the 3rd moon, while the sacrificial oblations
        were being prepared in the Hu Lu temple, a pan with oil would have caught
        fire, through the want of care on the part of the bonze, and that in a short
        time the flames would have consumed the paper pasted on the windows.}
}
{
}

\Verse{80}
{
    \zA{}
}
{
    \eA{Among the natives of this district bamboo fences and wooden partitions were
        in general use, and these too proved a source of calamity so ordained by
        fate (to consummate this decree). With promptness (the fire) extended to two
        buildings, then enveloped three, then dragged four (into ruin), and then
        spread to five houses, until the whole street was in a blaze, resembling the
        flames of a volcano. Though both the military and the people at once ran to
        the rescue, the fire had already assumed a serious hold, so that it was
        impossible for them to afford any effective assistance for its suppression.
        It blazed away straight through the night, before it was extinguished, and
        consumed, there is in fact no saying how many dwelling houses.}
}
{
}

\Verse{81}
{
    \zA{}
}
{
    \eA{Anyhow, pitiful to relate, the Chen house, situated as it was next door to
        the temple, was, at an early part of the evening, reduced to a heap of tiles
        and bricks; and nothing but the lives of that couple and several inmates of
        the family did not sustain any injuries. Shih-yin was in despair, but all he
        could do was to stamp his feet and heave deep sighs. After consulting with
        his wife, they betook themselves to a farm of theirs, where they took up
        their quarters temporarily. But as it happened that water had of late years
        been scarce, and no crops been reaped, robbers and thieves had sprung up
        like bees, and though the Government troops were bent upon their capture, it
        was anyhow difficult to settle down quietly on the farm.}
}
{
}

\Verse{82}
{
    \zA{}
}
{
    \eA{He therefore had no other resource than to convert, at a loss, the whole of
        his property into money, and to take his wife and two servant girls and come
        over for shelter to the house of his father-in-law. His father-in-law, Feng
        Su, by name, was a native of Ta Ju Chou. Although only a labourer, he was
        nevertheless in easy circumstances at home. When he on this occasion saw his
        son-in-law come to him in such distress, he forthwith felt at heart
        considerable displeasure.}
}
{
}

\Verse{83}
{
    \zA{}
}
{
    \eA{Fortunately Shih-yin had still in his possession the money derived from the
        unprofitable realization of his property, so that he produced and handed it
        to his father-in-law, commissioning him to purchase, whenever a suitable
        opportunity presented itself, a house and land as a provision for food and
        raiment against days to come. This Feng Su, however, only expended the half
        of the sum, and pocketed the other half, merely acquiring for him some
        fallow land and a dilapidated house. Shih-yin being, on the other hand, a
        man of books and with no experience in matters connected with business and
        with sowing and reaping, subsisted, by hook and by crook, for about a year
        or two, when he became more impoverished.}
}
{
}

\Verse{84}
{
    \zA{}
}
{
    \eA{In his presence, Feng Su would readily give vent to specious utterances,
        while, with others, and behind his back, he on the contrary expressed his
        indignation against his improvidence in his mode of living, and against his
        sole delight of eating and playing the lazy. Shih-yin, aware of the want of
        harmony with his father-in-law, could not help giving way, in his own heart,
        to feelings of regret and pain. In addition to this, the fright and vexation
        which he had undergone the year before, the anguish and suffering (he had
        had to endure), had already worked havoc (on his constitution); and being a
        man advanced in years, and assailed by the joint attack of poverty and
        disease, he at length gradually began to display symptoms of decline.}
}
{
}

\Verse{85}
{
    \zA{}
}
{
    \eA{Strange coincidence, as he, on this day, came leaning on his staff and with
        considerable strain, as far as the street for a little relaxation, he
        suddenly caught sight, approaching from the off side, of a Taoist priest
        with a crippled foot; his maniac appearance so repulsive, his shoes of
        straw, his dress all in tatters, muttering several sentiments to this
        effect:  All men spiritual life know to be good,  But fame to disregard they
        ne'er succeed! From old till now the statesmen where are they? Waste lie
        their graves, a heap of grass, extinct. All men spiritual life know to be
        good,  But to forget gold, silver, ill succeed!}
}
{
}

\Verse{86}
{
    \zA{}
}
{
    \eA{Through life they grudge their hoardings to be scant,  And when plenty has
        come, their eyelids close. All men spiritual life hold to be good,  Yet to
        forget wives, maids, they ne'er succeed! Who speak of grateful love while
        lives their lord,  And dead their lord, another they pursue. All men
        spiritual life know to be good,  But sons and grandsons to forget never
        succeed! From old till now of parents soft many,  But filial sons and
        grandsons who have seen? Shih-yin upon hearing these words, hastily came up
        to the priest, "What were you so glibly holding forth?" he inquired. "All I
        could hear were a lot of hao liao (excellent, finality.") "You may well have
        heard the two words 'hao liao,'" answered the Taoist with a smile, "but can
        you be said to have fathomed their meaning?}
}
{
}

\Verse{87}
{
    \zA{}
}
{
    \eA{You should know that all things in this world are excellent, when they have
        attained finality; when they have attained finality, they are excellent; but
        when they have not attained finality, they are not excellent; if they would
        be excellent, they should attain finality. My song is entitled
        Excellent-finality (hao liao)." Shih-yin was gifted with a natural
        perspicacity that enabled him, as soon as he heard these remarks, to grasp
        their spirit. "Wait a while," he therefore said smilingly; "let me unravel
        this excellent-finality song of yours; do you mind?" "Please by all means go
        on with the interpretation," urged the Taoist;}
}
{
}

\Verse{88}
{
    \zA{}
}
{
    \eA{whereupon Shih-yin proceeded in this strain:  Sordid rooms and vacant
        courts,  Replete in years gone by with beds where statesmen lay; Parched
        grass and withered banian trees,  Where once were halls for song and dance!
        Spiders' webs the carved pillars intertwine,  The green gauze now is also
        pasted on the straw windows! What about the cosmetic fresh concocted or the
        powder just scented; Why has the hair too on each temple become white like
        hoarfrost! Yesterday the tumulus of yellow earth buried the bleached bones,
        To-night under the red silk curtain reclines the couple! Gold fills the
        coffers, silver fills the boxes,  But in a twinkle, the beggars will all
        abuse you! While you deplore that the life of others is not long,  You
        forget that you yourself are approaching death!}
}
{
}

\Verse{89}
{
    \zA{}
}
{
    \eA{You educate your sons with all propriety,  But they may some day, 'tis hard
        to say become thieves; Though you choose (your fare and home) the fatted
        beam,  You may, who can say, fall into some place of easy virtue! Through
        your dislike of the gauze hat as mean,  You have come to be locked in a
        cangue; Yesterday, poor fellow, you felt cold in a tattered coat,  To-day,
        you despise the purple embroidered dress as long! Confusion reigns far and
        wide! you have just sung your part, I come on  the boards,  Instead of
        yours, you recognise another as your native land; What utter perversion! In
        one word, it comes to this we make wedding clothes for others! (We sow for
        others to reap.) The crazy limping Taoist clapped his hands. "Your
        interpretation is explicit," he remarked with a hearty laugh, "your
        interpretation is explicit!"}
}
{
}

\Verse{90}
{
    \zA{}
}
{
    \eA{Shih-yin promptly said nothing more than,--"Walk on;" and seizing the stole
        from the Taoist's shoulder, he flung it over his own. He did not, however,
        return home, but leisurely walked away, in company with the eccentric
        priest. The report of his disappearance was at once bruited abroad, and
        plunged the whole neighbourhood in commotion; and converted into a piece of
        news, it was circulated from mouth to mouth. Dame Feng, Shih-yin's wife,
        upon hearing the tidings, had such a fit of weeping that she hung between
        life and death; but her only alternative was to consult with her father, and
        to despatch servants on all sides to institute inquiries. No news was
        however received of him, and she had nothing else to do but to practise
        resignation, and to remain dependent upon the support of her parents for her
        subsistence.}
}
{
}

\Verse{91}
{
    \zA{}
}
{
    \eA{She had fortunately still by her side, to wait upon her, two servant girls,
        who had been with her in days gone by; and the three of them, mistress as
        well as servants, occupied themselves day and night with needlework, to
        assist her father in his daily expenses. This Feng Su had after all, in
        spite of his daily murmurings against his bad luck, no help but to submit to
        the inevitable. On a certain day, the elder servant girl of the Chen family
        was at the door purchasing thread, and while there, she of a sudden heard in
        the street shouts of runners clearing the way, and every one explain that
        the new magistrate had come to take up his office. The girl, as she peeped
        out from inside the door, perceived the lictors and policemen go by two by
        two;}
}
{
}

\Verse{92}
{
    \zA{}
}
{
    \eA{and when unexpectedly in a state chair, was carried past an official, in
        black hat and red coat, she was indeed quite taken aback. "The face of this
        officer would seem familiar," she argued within herself; "just as if I had
        seen him somewhere or other ere this." Shortly she entered the house, and
        banishing at once the occurrence from her mind, she did not give it a second
        thought. At night, however, while she was waiting to go to bed, she suddenly
        heard a sound like a rap at the door. A band of men boisterously cried out:
        "We are messengers, deputed by the worthy magistrate of this district, and
        come to summon one of you to an enquiry." Feng Su, upon hearing these words,
        fell into such a terrible consternation that his eyes stared wide and his
        mouth gaped.}
}
{
}

\Verse{93}
{
    \zA{}
}
{
    \eA{What calamity was impending is not as yet ascertained, but, reader, listen
        to the explanation contained in the next chapter.}
}
{
}
